% 第7章 現代ポートフォリオ理論(3)(資料7)
\chapter{現代ポートフォリオ理論 (3)：CAPMと市場モデル}

本章では，前章で導入した接点ポートフォリオの特別な性質を導き，それをもとに
市場均衡における資産価格の関係式である \textbf{CAPM}(Capital Asset Pricing
Model，資本資産価格モデル)を導出する。さらに，均衡理論とは別の実務的な枠組みで
ある\textbf{市場モデル}・\textbf{ファクターモデル}と，その応用である
インデックス運用・スタイル運用を扱う。

\section{接点ポートフォリオの性質}

\subsection{一般に成り立つ関係式}

まず，\emph{任意の}ポートフォリオと，そこに含まれる資産の収益率の間で一般に
成り立つ関係式を4つ示す。ポートフォリオの収益率を $R_P$，資産Aの収益率を
$R_A$，資産Aのウェイト(価値で測った保有比率)を $x_A$ とする。

\begin{enumerate}
\item \textbf{ポートフォリオの分散と共分散 (1)}：ポートフォリオの収益率 $R_P$ と
資産Aの収益率 $R_A$ の共分散の分解式
\begin{equation}
C[R_A, R_P] = x_A V[R_A] + \sum_{i\neq A} x_i\, C[R_A, R_i]
= \sum_i x_i\, C[R_A, R_i].
\label{eq:rel1}
\end{equation}
\item \textbf{ポートフォリオの分散と共分散 (2)}：ポートフォリオの収益率の分散は，
含まれている各資産とポートフォリオの収益率の共分散の(投資比率を重みとした)
加重平均
\begin{equation}
V[R_P] = \sum_i x_i\, C[R_i, R_P].
\label{eq:rel2}
\end{equation}
\item \textbf{ポートフォリオの標準偏差と相関 (1)}：ポートフォリオのリスク
($=$収益率の標準偏差)は，相関係数により調整された各資産のリスクの
(投資比率を重みとした)加重平均
\begin{equation}
\sigma[R_P] = \sum_i x_i\, \rho[R_i, R_P]\,\sigma[R_i].
\label{eq:rel3}
\end{equation}
\item \textbf{ポートフォリオの標準偏差と相関 (2)}：リスクの投資比率に対する感応度
\begin{equation}
\frac{\partial\sigma[R_P]}{\partial x_i} = \rho[R_i, R_P]\,\sigma[R_i].
\label{eq:rel4}
\end{equation}
すなわち，ポートフォリオに資産Aを僅かに($\Delta$)増やすと，ポートフォリオの
リスクは $\rho[R_A,R_P]\sigma[R_A]\Delta$ だけ増加する。
\end{enumerate}

\begin{proof}[証明]
\textbf{1.}共分散の定義から，$R_P = \sum_i x_i R_i$ を代入して
\begin{align*}
C[R_A, R_P] &= \E\bigl[(R_A - \E[R_A])(R_P - \E[R_P])\bigr]\\
&= \E\Bigl[(R_A-\E[R_A])\sum_i x_i(R_i - \E[R_i])\Bigr]\\
&= x_A\E\bigl[(R_A-\E[R_A])^2\bigr]
+ \sum_{i\neq A}x_i\E\bigl[(R_A-\E[R_A])(R_i-\E[R_i])\bigr]\\
&= x_A V[R_A] + \sum_{i\neq A}x_i C[R_A,R_i].
\end{align*}

\textbf{2.}1の結果を使うと，
\[
\sum_i x_i\,C[R_i,R_P]
= \sum_i x_i\Bigl(x_i V[R_i] + \sum_{j\neq i}x_j C[R_i,R_j]\Bigr)
= V[R_P].
\]

\textbf{3.}2の結果から，
\[
\sigma[R_P] = \frac{V[R_P]}{\sigma[R_P]}
= \sum_i x_i\,\frac{C[R_i,R_P]}{\sigma[R_P]}
= \sum_i x_i\,\rho[R_i,R_P]\,\sigma[R_i].
\]
最後の等号では $C[R_i,R_P] = \rho[R_i,R_P]\sigma[R_i]\sigma[R_P]$ を使った。

\textbf{4.}3の式を $x_i$ で偏微分すればよい(別解は本章の演習問題1)。
\end{proof}

\begin{remark}[リスク分解の意味]
ここまでの関係式は，任意のポートフォリオと，そこに含まれる資産の間の関係式で
あり，リスク管理上ある程度深い意味がある。ポートフォリオの\emph{リターン}は
常に各資産の寄与に分解できる(期待値の線形性が成り立つから)。一方，ポート
フォリオ全体の\emph{リスク}は，各資産の寄与に分解できるとは限らず，ある条件を
満たすことが必要である。しかし，標準偏差をリスクとして選んだ場合には，式
\eqref{eq:rel3}のとおりポートフォリオ全体のリスクを各資産の寄与に分解できる
(標準偏差が分解可能な条件を満たすためである)。
\end{remark}

\vspace{18pt}

\subsection{接点ポートフォリオの分散・共分散の性質}

接点ポートフォリオの値を添字 $T$ で示すと，次式が成り立つ。

\begin{theorem}[接点ポートフォリオの性質]
$\mu_T$ と $\mu_i$ を接点ポートフォリオと資産 $i$ の期待収益率とすると，
\begin{equation}
\frac{\mu_T - r}{\sigma[R_T]}
= \frac{\mu_i - r}{\rho[R_i,R_T]\,\sigma[R_i]},
\qquad i = 1,\dots,N.
\label{eq:tangent-property}
\end{equation}
すなわち，各資産の「(期待超過収益率)/(ポートフォリオのリスクに対する感応度
\eqref{eq:rel4})」は一定で，効率的フロンティアの傾きに等しい。
\end{theorem}

\begin{proof}
Tobinの問題の解法(前章の補足)を利用する。接点ポートフォリオも一つの最小分散
ポートフォリオなので，Lagrange法の1階条件より
\[
\sum_{k=1}^{N} C[R_i,R_k]\,x_k - \lambda_1\mu_i - \lambda_2 = 0,
\qquad i=1,\dots,N,
\qquad\text{かつ}\quad
-\lambda_1 r - \lambda_2 = 0
\]
が成り立つ。後者($\lambda_2 = -r\lambda_1$)を前者に代入すると，
式\eqref{eq:rel1}を使って
\begin{equation}
C[R_i, R_T] = \sum_{k=1}^{N} C[R_i,R_k]\,x_k = \lambda_1(\mu_i - r),
\qquad i = 1,\dots,N.
\label{eq:tangent-foc}
\end{equation}
よって
\begin{equation}
\frac{\sigma[R_T]}{\lambda_1}
= \frac{(\mu_i - r)\,\sigma[R_T]}{C[R_i,R_T]}
= \frac{(\mu_i - r)\,\sigma[R_T]}{\rho[R_i,R_T]\sigma[R_i]\sigma[R_T]}
= \frac{\mu_i - r}{\rho[R_i,R_T]\,\sigma[R_i]},
\qquad i=1,\dots,N
\label{eq:tangent-star}
\end{equation}
が成り立つ。さらに，\eqref{eq:tangent-foc}の両辺に $x_i$ を掛けて $i$ について
和をとると，式\eqref{eq:rel2}より左辺は
$\sum_i x_i C[R_i,R_T] = V[R_T] = \sigma^2[R_T]$，右辺は
$\lambda_1\sum_i x_i(\mu_i - r) = \lambda_1(\mu_T - r)$ となるので，
\[
\sigma^2[R_T] = \lambda_1(\mu_T - r)
\qquad\therefore\quad
\frac{\sigma[R_T]}{\lambda_1} = \frac{\mu_T - r}{\sigma[R_T]}.
\]
両者は等しくなければならないので，\eqref{eq:tangent-star}とあわせて
\eqref{eq:tangent-property}を得る。
\end{proof}

いま証明した式\eqref{eq:tangent-property}は，接点ポートフォリオの「リスクと
超過リターンの比」が，各資産の「(相関係数で調整後の)リスクと超過リターンの比」
に等しいことを示している。

\begin{remark}
この証明では「最小分散ポートフォリオであること」を使っているだけなので，
式\eqref{eq:tangent-property}は接点ポートフォリオ $T$ 以外の最小分散ポート
フォリオでも成立する。
\end{remark}

\vspace{18pt}

\section{CAPM(資本資産価格モデル)}

CAPMは，\textbf{均衡時における資産価格間の関係式}である。

\subsection{仮定}

以下では，市場は\textbf{均衡状態}(需要と供給が完全に釣り合った状態。ここでは
全ての資産は必ず誰かに保有されていて，余剰資産はないと考える)とする。

\begin{keypoint}[CAPMの仮定]
\begin{enumerate}
\item すべての投資家はリスク回避的であるとして，平均分散モデルによる最適投資を
考える。
\item すべての投資家は，証券のリスクとリターンと共分散について同一の期待を
持っている($=$同じ情報を持っている)。
\item すべての投資家は，同一の利子率 $r$ で任意の金額の貸借ができる
($=$利回り $r$ の無リスク資産を任意の量取引可能)。
\item 市場は完全である。
\end{enumerate}
\end{keypoint}

ここで\textbf{完全市場}とは，以下を満たす市場のことである。
\begin{enumerate}
\item[①] 投資家はプライステイカーである(個々の投資家の取引は市場価格に影響を
与えない)。
\item[②] 証券は無限分割可能(任意の量の取引が可能)である。
\item[③] 取引費用や税金は存在しない。
\item[④] 投資家は自由に空売りができる(負の保有量も可能)。
\end{enumerate}

\subsection{均衡状態下の接点ポートフォリオ=市場ポートフォリオ}

仮定1， 2より，すべての投資家の効率的フロンティアは共通である。仮定1--4より，
Tobinの問題の結果を使うことができて，\textbf{全ての投資家の接点ポートフォリオが
同じになる}。
\begin{itemize}
\item Tobinのポートフォリオ選択問題では，全投資家の最適ポートフォリオは接点
ポートフォリオと無リスク資産の組み合わせで構成される。
\item どの投資家の最適ポートフォリオでも，リスク性資産の部分だけを考えて構成
比率を見ると，接点ポートフォリオの構成比率に等しい(リスク性資産と無リスク資産の
構成比率は投資家のリスク選好による)。
\end{itemize}
したがって，\textbf{接点ポートフォリオは(リスク性資産の)市場ポートフォリオ}で
ある。ここで\textbf{市場ポートフォリオ}とは，市場にある全ての資産を，市場全体に
おける当該資産の構成比率と同じ構成比率で持つポートフォリオのことである。

\subsection{資本市場線 (CML)}

均衡状態における投資家の(効率的)ポートフォリオ P は
\begin{equation}
\mu_P = r + \frac{\mu_M - r}{\sigma_M}\,\sigma_P
\label{eq:cml}
\end{equation}
を満たす。ただし，$(\sigma_P,\mu_P)$ は投資家のポートフォリオPのリスクと
リターン，$(\sigma_M,\mu_M)$ は市場ポートフォリオMのリスクとリターンである。

\begin{definition}[資本市場線とリスクの市場価格]
式\eqref{eq:cml}を\textbf{資本市場線} (Capital Market Line， CML) といい，
その傾き
\begin{equation}
\frac{\mu_M - r}{\sigma_M}
\end{equation}
を\textbf{リスクの市場価格} (market price of risk) という。CMLは
$(\sigma,\mu)$ 平面上の半直線である。
\end{definition}

これがCAPMの一つの大きな成果である。

\subsection{CAPMの関係式}

均衡における資産間のリスク・リターンの関係式を整理する。任意のポートフォリオで
成立する式はもちろん成立する。また，市場ポートフォリオMは接点ポートフォリオTと
一致するので，最小分散ポートフォリオで成り立つ式も成立する。具体的には，
リスクに関して(式\eqref{eq:rel3}，\eqref{eq:rel4}より)
\begin{equation}
\sigma[R_M] = \sum_i x_i\,\rho[R_i,R_M]\,\sigma[R_i]
= \sum_i x_i\,\frac{\partial\sigma[R_M]}{\partial x_i},
\label{eq:capm3}
\end{equation}
リスクとリターンの関係として(式\eqref{eq:tangent-property}で $T\to M$ として)
\begin{equation}
\frac{\mu_M - r}{\sigma[R_M]}
= \frac{\mu_i - r}{\rho[R_i,R_M]\,\sigma[R_i]},
\qquad i = 1,\dots,N
\label{eq:capm4}
\end{equation}
が成立する。

\paragraph{CAPMの解釈と個別資産のリスク。}
ポートフォリオMに含まれる資産 $i$ のリスクを
$\rho[R_i,R_M]\,\sigma[R_i]$ とみなすと，式\eqref{eq:capm4}は，個々の資産に
おける「期待超過収益率/リスク」が一定であることを示している。資産 $i$ を
ポートフォリオM自身とみなすと，
$\rho[R_M,R_M]\sigma[R_M] = \sigma[R_M]$ なので，Mの「期待超過収益率/リスク」も
式\eqref{eq:capm4}で表現できる。さらに，これらの関係式は，市場ポートフォリオMの
部分集合(部分ポートフォリオという)を資産 $i$ とみなしても成立する。

\subsection{証券市場線とベータ}

式\eqref{eq:capm4}より，
\begin{equation}
\mu_i - r = \beta_i(\mu_M - r),
\qquad
\beta_i = \frac{\sigma[R_i]\,\rho[R_i,R_M]}{\sigma[R_M]}
\label{eq:sml}
\end{equation}
が成り立つ。

\begin{definition}[証券市場線とベータ]
式\eqref{eq:sml}を\textbf{証券市場線} (Security Market Line， SML) と呼び，
\begin{equation}
\beta_i = \frac{\rho[R_i,R_M]\,\sigma[R_i]}{\sigma[R_M]}
= \frac{C[R_i,R_M]}{V[R_M]}
\end{equation}
を証券 $i$ の\textbf{ベータ}($\beta$)と呼ぶ。
\end{definition}

これもCAPMの一つの大きな成果である。$(\beta,\mu)$ 平面上でSMLを描くと，切片
$r$，傾き $\mu_M - r$ の半直線であり，$\beta=1$ のとき $\mu = \mu_M$(市場
ポートフォリオ)を通る。$\beta$ が示すリスクを\textbf{ベータ・リスク}という。
この直線は，証券市場における「ハイリスク・ハイリターン，ローリスク・ロー
リターン」の関係を示している。

\begin{remark}
SMLは$(\beta,\mu)$平面上の半直線，CMLは$(\sigma,\mu)$平面上の半直線である。
混同しないこと。
\end{remark}

式\eqref{eq:sml}はまた
\begin{equation}
\mu_i - r = \rho[R_i,R_M]\,\sigma[R_i] \times \frac{\mu_M - r}{\sigma[R_M]}
\end{equation}
とも書ける。すなわち，
\[
(\text{資産 } i \text{ の期待超過収益率})
= (\text{資産 } i \text{ の(相関で調整後の)リスク})
\times(\text{リスクの市場価格})
\]
である。均衡時には，すべての証券は証券市場線上に並ぶ。$\beta_i$ は，市場
ポートフォリオを基準としたリスク尺度でもある。

\subsection{CAPMの2つの定理}

\begin{theorem}[CAPM第一定理]
無リスク資産が存在するとき，市場が均衡状態ならば，市場ポートフォリオは接点
ポートフォリオと一致する。ゆえに，市場ポートフォリオは効率的ポートフォリオで
ある。
\end{theorem}

\begin{theorem}[CAPM第二定理]
無リスク資産が存在するとき，市場が均衡状態ならば，証券市場線の式\eqref{eq:sml}が
成立する。
\end{theorem}

\vspace{18pt}

\subsection{ベータの推定}

$\beta$ は，$(r, R_M, R_i)$ の時系列データから回帰分析を使って推定可能である。
$(R_M - r,\ R_i - r)$ の時系列データをとり，
\begin{equation}
(R_i - r) = \alpha_i + \beta_i(R_M - r) + \varepsilon_i
\end{equation}
と仮定して回帰分析を行う。このようにして求めた $\beta$ を\textbf{ヒストリカル・
ベータ}といい，$\alpha$ を\textbf{ヒストリカル・アルファ}という。$\alpha$ は
\textbf{Jensen(ジェンセン)のアルファ}とも呼ばれる。

\subsection{CAPMと割高・割安分析}

CAPM理論は「均衡時における」証券価格の話である。概念的には，
\begin{itemize}
\item 過去の実績が証券市場線より\textbf{上}にプロットされる資産は\textbf{割安}
(価格が安いために実績リターンが理論値より高い)，
\item 証券市場線より\textbf{下}にプロットされる資産は\textbf{割高}，
\end{itemize}
と考えられる。ただし，CAPMの理論式は実際の市場の中で厳密に成立しているわけでは
ない。個々の資産の期待収益率は，市場ポートフォリオのリターン以外のことにも
依存する。

\subsection{ゼロベータCAPMへの一般化}

CAPM第二定理は，以下のように一般化できる。
\begin{itemize}
\item 無リスク資産がなくても市場が均衡状態ならば，式\eqref{eq:sml}と同じ形の式が
成立する。ただし，ゼロ・ベータ・ポートフォリオの収益率を無リスク金利 $r$ の
代替として使用する。
\item すなわち，無リスク資産の有無によらず，式\eqref{eq:sml}と同じ形の式は成立
する。
\end{itemize}

\begin{keypoint}[ゼロベータCAPM]
ある条件を満たすポートフォリオ(ゼロ・ベータ・ポートフォリオ)を無リスク資産の
代替として使うことで，通常のCAPMと同形式の関係式
\begin{equation}
\mu_i - \mu_Z = \beta_i(\mu_M - \mu_Z),\qquad
\beta_i = \frac{\sigma[R_i]\,\rho[R_i,R_M]}{\sigma[R_M]}
\end{equation}
を導出できる。ここで，添字 $Z$ はゼロ・ベータ・ポートフォリオにおける値である。
証明や詳細(ゼロ・ベータ・ポートフォリオとは何か)は第5章の補足
(Markowitz-5.pdf)を参照。理論上は，基準となるポートフォリオとして最小分散
境界上の任意のポートフォリオ(大域的最小分散ポートフォリオを除く)を選ぶことが
できる。
\end{keypoint}

\section{市場モデルとファクター}

ここで頭を切り替える必要がある。\textbf{本節では均衡から離れる}。

\subsection{市場モデル，ファクターモデル}

\begin{definition}
\textbf{市場効果}とは，株価指数(日経平均，TOPIXなど)の変動に連動して株価が
変動する現象のこと。このような銘柄はかなり多い。
\end{definition}

ある要因の変動に応じて各銘柄の株価が変動するモデルを考える。
\begin{itemize}
\item \textbf{ファクターモデル}：変動要因が一つまたは複数あるモデル。変動要因は
モデル立案者が自分で選択する。
\item \textbf{市場モデル}：変動要因として株価指数を選んだモデル(1ファクター
モデル)。ファクターモデルの一例。
\end{itemize}

\subsection{市場モデルの定式化}

個別証券 $j$ の収益率 $R_j$ と市場ポートフォリオの収益率 $R_M$ の間には，次の
関係式が成り立つと\emph{仮定}する：
\begin{equation}
R_j = \alpha_j + \beta_j R_M + \varepsilon_j,\qquad j = 1,\dots,n.
\label{eq:market-model}
\end{equation}
ただし，$\alpha_j, \beta_j$ は定数で，
\begin{align*}
&\E[\varepsilon_j] = 0,\quad V[\varepsilon_j] = \sigma_{\varepsilon,j}^2,
\quad j=1,\dots,n,\\
&C[\varepsilon_j,\varepsilon_k] = 0\ (j\neq k),\qquad
C[\varepsilon_j, R_M] = 0,\quad j=1,\dots,n
\end{align*}
とする($\sigma_{\varepsilon,j}^2$ も定数)。$\varepsilon_j$ は，市場ポート
フォリオの変動では説明できない証券固有の変動を表す確率変動項である。また
$\mu_M = \E[R_M]$， $\sigma_M^2 = V[R_M]$， $\mu_j = \E[R_j]$，
$\sigma_j^2 = V[R_j]$ と書く。要するに，回帰分析である。

このモデルから，資産 $j$ のリターン(期待収益率)と分散は
\begin{align}
\mu_j &= \E[R_j] = \alpha_j + \beta_j \E[R_M],\qquad j=1,\dots,n,\\
\sigma_j^2 &= V[R_j] = \beta_j^2\sigma_M^2 + \sigma_{\varepsilon,j}^2,
\qquad j = 1,\dots,n
\label{eq:mm-var}
\end{align}
となる。分散の分解\eqref{eq:mm-var}において，
\begin{itemize}
\item $\beta_j^2 V[R_M]$:\textbf{システマティックリスク}(の寄与)。市場
ポートフォリオの収益率に連動する成分の寄与。
\item $\sigma_{\varepsilon,j}^2$:\textbf{非システマティックリスク}(の寄与)。
資産固有の成分の寄与。
\item $\sigma_j^2$:\textbf{総リスク} (total risk)。
\end{itemize}

\subsection{回帰分析によるパラメータ推定}

時系列データ $(R_j, R_M)$ を回帰分析して
$\alpha_j, \beta_j, \sigma_{\varepsilon,j}$ を求める。回帰方程式は
$R_j = \alpha_j + \beta_j R_M$($j=1,\dots,n$)であり，最小二乗法で推定可能で
ある。回帰係数は
\begin{equation}
\beta_j = \frac{C[R_M,R_j]}{V[R_M]} = \frac{\sigma_{Mj}}{\sigma_M^2},\qquad
\alpha_j = \E[R_j] - \beta_j\E[R_M]
\end{equation}
となる。また，決定係数(総リスク $\sigma_j^2$ に占めるシステマティック・リスクの
割合)は
\begin{equation}
\rho^2 = \rho^2[R_M,R_j]
= \frac{\beta_j^2\sigma_M^2}{\sigma_j^2}
= 1 - \frac{\sigma_{\varepsilon,j}^2}{\sigma_j^2}
\end{equation}
である。

\subsection{CAPMと市場モデルは全く違う}

\begin{keypoint}[CAPMと市場モデルの違い]
\begin{itemize}
\item \textbf{CAPM}は，厳しい仮定の下で導出された，\emph{均衡時に成立する式}で
ある。均衡時の合理的なポートフォリオのリターンとリスクの間の関係式(資本市場線
$\mu_P - r = \frac{\mu_M - r}{\sigma_M}\sigma_P$)と，均衡時の個別証券の
リターンとリスクの間の関係式(証券市場線
$\mu_A - r = \beta_A(\mu_M - r)$， $\beta_A = \rho_{AM}\sigma_A/\sigma_M$)を
与える。
\item 一方，\textbf{市場モデル}は\emph{単なる仮定として立てられた式}に過ぎない。
証券市場線と同形の式自体が仮定であり，議論の出発点になっている。ただし，市場
参加者の感覚にあう式でもある。設定は自由で応用も幅広く，投資技法として広く活用
されている。
\end{itemize}
通常，個別資産の決定係数は低いが，ある程度分散されたポートフォリオでは決定係数が
高くなる(非システマティックリスクが低下するため)。
\end{keypoint}

\subsection{CAPMの反例とFama-Frenchの3ファクターモデル}

実証研究により，株式市場には\textbf{アノマリー}が発見されている。すなわち，
投資収益率は市場ファクター以外にも依存する。

その代表的な定式化が \textbf{Fama-Frenchの3ファクターモデル}である：
\begin{equation}
R_{j,t} - r = \alpha_j + \beta_j^{MKT}(R_{M,t} - r)
+ \beta_j^{SMB}\,SMB_t + \beta_j^{HML}\,HML_t + \varepsilon_{j,t}.
\end{equation}
ここで，
\begin{itemize}
\item $r$：無リスク金利，$\alpha_j,\beta_j^{MKT},\beta_j^{SMB},\beta_j^{HML}$：
株価 $j$ ごとの定数，
\item $R_{j,t}$：時点 $t$ の株価 $j$ の収益率，
\item $R_{M,t}$：時点 $t$ の市場ポートフォリオの収益率，
\item $SMB_t$：時点 $t$ の時価総額リスクファクター(時価総額で下位50\%ポートの
株価収益率 $-$ 上位50\%ポートの株価収益率)，
\item $HML_t$：時点 $t$ の簿価/時価比率リスクファクター(簿価時価比率で上位
30\%ポートの株価収益率 $-$ 下位30\%ポートの株価収益率)，
\end{itemize}
であり，各々を\textbf{マーケット・ファクター}，\textbf{サイズ・ファクター}，
\textbf{バリュー・ファクター}と呼ぶ。このモデルは市場を説明する有効なモデルと
して広く認められている。なお，このモデルは APT (Arbitrage Pricing Theory) の例
としても語られるが，本講義ではAPTには触れないので，実務的な例として挙げる。

\subsection{ファクター投資戦略}

\begin{itemize}
\item \textbf{ファクター・ティルト(orファクター投資)戦略}：あるファクターへの
感応度が高いポートフォリオを構築し，そのファクターから超過収益を挙げようとする
戦略。Fama-French 3ファクターモデル(マーケット，サイズ，バリューに投資)が
代表例。他にも，ボラティリティ・ファクター，リスクプレミアム・ファクター
(低格付債と国債のスプレッド)，タームプレミアム・ファクター(長期債と短期債の
スプレッド)など，ファクターは多数ありうる。
\item \textbf{投資スタイル}：アクティブとパッシブ(後述)，バリュー(割安株)型，
グロース(成長株)型，など。
\item \textbf{スマートベータ}：特定のファクターに関するインデックス(スマート
ベータ指数)をベンチマークとする投資手法。バリュー，高配当，モメンタム，
低リスク，などがある。
\end{itemize}

\section{インデックス運用，スタイル運用}

\subsection{インデックス・ファンド}

均衡時の市場ポートフォリオは最小分散ポートフォリオである。しかし，自分で市場
ポートフォリオを維持するのは困難である。
\begin{itemize}
\item 全銘柄に投資しなければならない。投資額も取引コストも相当かかる。
\item 市場の変動にあわせてポートフォリオを随時組み替えねばならない。
\end{itemize}

その代替として，\textbf{インデックス・ファンド}を考える。これは，インデックスの
収益率 $R_M$ とほぼ同じ収益率を\emph{少数銘柄}で実現するポートフォリオである。
こうしたファンドを用いた運用手法を\textbf{インデックス運用}という。

参考として，次の用語もある。
\begin{itemize}
\item \textbf{パッシブ運用}：インデックスに連動した収益を求める運用。
\item \textbf{アクティブ運用}：インデックスを超える収益を積極的に求める運用。
\end{itemize}

\subsection{インデックス・ファンドの定式化}

できるだけ少ない資産，取引しやすい(流動性の高い)資産でインデックスファンドを
作るにはどうすればよいか。市場モデルにもとづく定式化の例を示す。

$x_j$ を資産 $j$ への投資比率とすると，ファンドの投資収益率は
\begin{equation}
R_P = \sum_{j=1}^{m} x_j(\alpha_j + \beta_j R_M + \varepsilon_j)
\end{equation}
であり，このファンドの収益率の分散は
\begin{equation}
\sigma_p^2 = \Bigl(\sum_{j=1}^{m}x_j\beta_j\Bigr)^2\sigma_M^2
+ \sum_{j=1}^{m}x_j^2\sigma_{\varepsilon,j}^2
\end{equation}
となる。そこで，次の最適化問題の解が，求めるインデックス・ファンドの組成を
与える：
\begin{equation}
\min_{x_1,\dots,x_m}\ \sum_{j=1}^{m}x_j^2\,\sigma_{\varepsilon,j}^2
\qquad\text{s.t.}\quad
\sum_{j=1}^{m}x_j\beta_j = 1,\qquad \sum_{j=1}^{m}x_j = 1.
\label{eq:index-opt}
\end{equation}
この式の意味は次のとおりである。制約 $\sum_j x_j\beta_j = 1$ はファンド全体の
ベータを1(市場と同じ感応度)にすること，目的関数はインデックスで説明できない
固有変動(トラッキング・エラーの源泉)を最小にすることを表す。現実の投資では，
条件がさらに追加されることが多い(例えば $x_j\geq 0$(空売り禁止)など)。

\begin{example}[例題1：インデックスファンドの組成]
株式 A， B， C と市場ポートフォリオMの収益率を $R_i$， $i=A,B,C,M$ で表す。市場
モデル\eqref{eq:market-model}で過去のデータを回帰分析したところ，
\[
(\alpha_A,\alpha_B,\alpha_C) = (0,\ 0.01,\ 0.02),\qquad
(\beta_A,\beta_B,\beta_C) = (0.25,\ 0.25,\ -0.5)
\]
で，個別株のリスクは
$(\sigma_A,\sigma_B,\sigma_C) = (0.1,\ 0.1,\ 0.2)$ であった。市場ポートフォリオが
$(\sigma_M,\mu_M) = (0.2,\ 0.06)$ であるとき，市場ポートフォリオに模した
インデックスファンドの組成(株式A， B， Cの投資比率)を求めなさい。

\medskip
\noindent\textbf{解答。}
定式化\eqref{eq:index-opt}より，インデックスファンドを求めるには
$\sigma_{\varepsilon,j}^2$ が必要である。そこでまず相関係数を求めると，
\[
\rho_{iM} = \frac{\beta_i\,\sigma_M}{\sigma_i}
\]
より
\[
(\rho_{AM},\rho_{BM},\rho_{CM}) = (0.5,\ 0.5,\ -0.5).
\]
よって
\[
\sigma_{\varepsilon,j}^2 = \sigma_j^2(1 - \rho_{jM}^2)
\]
より
\[
(\sigma_{\varepsilon,A}^2,\sigma_{\varepsilon,B}^2,\sigma_{\varepsilon,C}^2)
= (0.0075,\ 0.0075,\ 0.03).
\]
これより(次の式では比率だけが重要なので簡略に書くと)，
\[
\min_{x_A,x_B,x_C}\ \frac{3}{4}x_A^2 + \frac{3}{4}x_B^2 + 3x_C^2
\qquad\text{s.t.}\quad
\frac{1}{4}x_A + \frac{1}{4}x_B - \frac{1}{2}x_C = 1,\qquad
x_A + x_B + x_C = 1
\]
の解がインデックスファンドの組成を与える。

これを解こう。第2の制約の $1/4$ 倍を第1の制約から引くと
\[
\frac{1}{4}(x_A+x_B) - \frac{1}{2}x_C = 1,\quad x_A+x_B = 1 - x_C
\ \Longrightarrow\
\frac{1}{4}(1-x_C) - \frac{1}{2}x_C = 1
\ \Longrightarrow\ x_C = -1.
\]
すると $x_A + x_B = 2$ であり，目的関数のうち $x_A,x_B$ に関わる部分
$\frac{3}{4}(x_A^2+x_B^2)$ は，和が一定のとき $x_A = x_B$ で最小になるので，
$x_A = x_B = 1$ を得る。よって，インデックスファンドの組成は
\[
(x_A,\ x_B,\ x_C) = (1,\ 1,\ -1)
\]
である(株式Cを空売りして，その資金も使って株式A・Bを購入する)。
(この解答は木島・後藤・鈴木テキストの例5.4(105ページ)にも対応する。)
\end{example}

\vspace{18pt}

\subsection{現代ポートフォリオ理論の最後に}

\begin{keypoint}[まとめの注意]
\begin{itemize}
\item CAPMのポイントをグラフで表現した，(1) CML(資本市場線)，(2) SML
(証券市場線)の違いに注意すること。
\item CAPMは市場が\emph{均衡状態}にあるときの話である。
\item CAPMのSMLの式と，市場モデルの式の\emph{意味}の違いに注意すること。
\end{itemize}
\end{keypoint}

\section*{付録：CAPMの方程式の別の証明}

多数の証券が存在する場合の直接的な証明を示す(内容は第5章補足のゼロベータCAPMの
節と同じ論法である)。

市場ポートフォリオMと，任意のリスク証券A($(\sigma_A,\mu_A)$)の組み合わせに
よるポートフォリオを考える。半直線FM(F：無リスク資産)は効率的フロンティアで
あるから，この線より上に投資機会集合は存在しない。そのため，MとAを組み合わせた
ポートフォリオが描く曲線AMは，点Mにおいて半直線FMと\emph{接する}(曲線AMの点Mに
おける接線は半直線FMと一致する)。もし接せずに交わるならば，半直線FMより上に
投資機会が存在することになって矛盾するからである。このことを数式で表現する。

\paragraph{手順1：曲線AMを求める。}
リスク証券Aへの投資比率を $w$ とすると M : A $= (1-w) : w$ であり，ポート
フォリオの収益率・リターン・リスクは
\begin{align}
R_P &= (1-w)R_M + wR_A,\\
\mu_P &= \E[R_P] = (1-w)\mu_M + w\mu_A,\\
\sigma_P^2 &= V[R_P] = (1-w)^2\sigma_M^2 + 2w(1-w)\sigma_{AM} + w^2\sigma_A^2,
\end{align}
ただし $\sigma_{AM} = \Cov[R_M,R_A] = \rho_{AM}\sigma_M\sigma_A$ である。よって
\begin{equation}
\sigma_P = \sqrt{(1-w)^2\sigma_M^2 + 2w(1-w)\sigma_{AM} + w^2\sigma_A^2}.
\end{equation}

\paragraph{手順2：接線の傾きを求める。}
接線の傾きは
\[
\frac{\dd\mu_P}{\dd\sigma_P} = \frac{\dd\mu_P/\dd w}{\dd\sigma_P/\dd w}
\]
と書ける。分子は
\[
\frac{\dd\mu_P}{\dd w} = \mu_A - \mu_M,
\]
分母は
\[
\frac{\dd\sigma_P^2}{\dd w} = 2\sigma_P\frac{\dd\sigma_P}{\dd w}
= -2(1-w)\sigma_M^2 + 2(1-2w)\sigma_{AM} + 2w\sigma_A^2
\]
より
\[
\frac{\dd\sigma_P}{\dd w}
= \frac{-(1-w)\sigma_M^2 + (1-2w)\sigma_{AM} + w\sigma_A^2}{\sigma_P}
\]
なので，
\begin{equation}
\frac{\dd\mu_P}{\dd\sigma_P}
= \frac{\sigma_P(\mu_A - \mu_M)}
{-(1-w)\sigma_M^2 + (1-2w)\sigma_{AM} + w\sigma_A^2}.
\end{equation}

\paragraph{手順3：点Mでの一致条件。}
点Mでは市場ポートフォリオ100\%なので $w=0$ を代入すると
($\sigma_P|_{w=0} = \sigma_M$)，
\begin{equation}
\left.\frac{\dd\mu_P}{\dd\sigma_P}\right|_{w=0}
= \frac{\sigma_M(\mu_A - \mu_M)}{\sigma_{AM} - \sigma_M^2}.
\end{equation}
これが半直線FMの傾き $(\mu_M - r_F)/\sigma_M$ に等しいので，
\[
\frac{\sigma_M(\mu_A-\mu_M)}{\sigma_{AM}-\sigma_M^2}
= \frac{\mu_M - r_F}{\sigma_M}.
\]
よって，
\begin{equation}
\mu_A = \mu_M + \frac{\mu_M - r_F}{\sigma_M^2}(\sigma_{AM} - \sigma_M^2)
= r_F + \frac{\sigma_{AM}}{\sigma_M^2}(\mu_M - r_F)
= r_F + \frac{\rho_{AM}\sigma_A}{\sigma_M}(\mu_M - r_F)
\end{equation}
が得られる。これは証券市場線の式\eqref{eq:sml}そのものである。

% ---------------- 演習問題 ----------------
\section{演習問題}

\begin{exercise}[問題1]
次の式
\[
\frac{\partial\sigma[R_P]}{\partial x_i} = \rho[R_i,R_P]\,\sigma[R_i]
\]
を，式\eqref{eq:rel3}を $x_i$ で偏微分する方法\emph{以外}で証明しなさい。
\end{exercise}

\begin{answerbox}
\textbf{考え方。}標準偏差の代わりに\emph{分散}を偏微分してみる。分散
$\sigma^2[R_P]$ は投資比率の二次形式なので偏微分は容易であり，一方
$\sigma^2[R_P] = (\sigma[R_P])^2$ の合成関数の微分と比較すれば，標準偏差の
偏微分が取り出せる。

\medskip
\textbf{解答。}
まず，ポートフォリオの収益率の分散を形式的に偏微分すると，
\begin{equation}
\frac{\partial\sigma^2[R_P]}{\partial x_i}
= 2\sigma[R_P]\,\frac{\partial\sigma[R_P]}{\partial x_i}.
\tag{1}
\end{equation}
一方，$R_P = \sum_j x_j R_j$ より，ポートフォリオの収益率の分散は
\[
\sigma^2[R_P] = \sum_{i,j} C[R_i,R_j]\,x_i x_j
\]
なので，
\begin{equation}
\frac{\partial\sigma^2[R_P]}{\partial x_i}
= \frac{\partial}{\partial x_i}\sum_{i,j}C[R_i,R_j]x_ix_j
= 2\sum_{j} C[R_i,R_j]\,x_j = 2\,C[R_i,R_P].
\tag{2}
\end{equation}
(最後の等号は関係式\eqref{eq:rel1}による。)

(1)と(2)より，
\[
\frac{\partial\sigma[R_P]}{\partial x_i}
= \frac{C[R_i,R_P]}{\sigma[R_P]}
= \frac{\rho[R_i,R_P]\,\sigma[R_i]\,\sigma[R_P]}{\sigma[R_P]}
= \rho[R_i,R_P]\,\sigma[R_i].
\qquad\blacksquare
\]
\end{answerbox}

\begin{exercise}[問題2]
市場ポートフォリオと資産Aの収益率をそれぞれ $R_M$， $R_A$，その期待値を
$\mu_M = 5\%$，$\mu_A$(未知)，標準偏差を $\sigma_M = 20\%$，
$\sigma_A = 15\%$，共分散を $\sigma_{MA} = 0.015$，無リスク金利を $r = 4\%$ と
する。
\begin{enumerate}
\item[(1)] $R_M$ と $R_A$ の相関係数 $\rho_{MA}$ を求めなさい。
\item[(2)] 市場ポートフォリオのリスクの市場価格を求めなさい。
\item[(3)] 株式Aのベータを求めなさい。
\item[(4)] CAPMが成立するとき，$\mu_A$ を求めなさい。
\item[(5)] $(\sigma,\mu)$ 平面上で資本市場線を描きなさい。また，資産Aの位置を
プロットしなさい。
\item[(6)] ある期間のデータ分析では $\mu_A = 4.5\%$ であった。証券市場線を描き，
この期間の資産Aは割安か割高かを述べなさい。
\end{enumerate}
\end{exercise}

\begin{answerbox}
\textbf{考え方。}(1)〜(4)は定義式への代入である。(5)(6)では，CML(横軸
$\sigma$)とSML(横軸 $\beta$)のどちらの平面で考えているかを意識することが
重要である。

\medskip
\textbf{(1)}相関係数は
\[
\rho_{MA} = \frac{\sigma_{MA}}{\sigma_M\sigma_A}
= \frac{0.015}{0.2\times 0.15} = 0.5.
\]

\textbf{(2)}リスクの市場価格は
\[
\lambda = \frac{\mu_M - r}{\sigma_M} = \frac{0.05 - 0.04}{0.2} = 0.05.
\]

\textbf{(3)}株式Aのベータは
\[
\beta_A = \frac{\rho_{MA}\,\sigma_A}{\sigma_M} = \frac{0.5\times 0.15}{0.2}
= 0.375,
\qquad\text{あるいは}\quad
\beta_A = \frac{\sigma_{MA}}{\sigma_M^2} = \frac{0.015}{0.2\times 0.2} = 0.375.
\]

\textbf{(4)}CAPM(証券市場線)が成立するとき，
\[
\mu_A = r + \beta_A(\mu_M - r) = 0.04 + 0.375\times(0.05-0.04)
= 0.04375 = 4.375\%.
\]

\textbf{(5)}資本市場線は，$(\sigma,\mu)$ 平面上で無リスク資産 $(0, 0.04)$ から
市場ポートフォリオ $(0.2, 0.05)$ を通る半直線
\[
\mu = 0.04 + 0.05\,\sigma
\]
である。資産Aは点 $(\sigma_A,\mu_A) = (0.15,\ 0.04375)$ にプロットされる。
$\sigma = 0.15$ における資本市場線上の点は $(0.15,\ 0.0475)$ なので，資産Aは
資本市場線より\emph{下}に位置する($\rho_{MA}<1$ なので効率的ポートフォリオでは
ない。これは自然なことである)。

\textbf{(6)}証券市場線は，$(\beta,\mu)$ 平面上で無リスク資産 $(0, 0.04)$ から
市場ポートフォリオ $(1, 0.05)$ を通る半直線
\[
\mu = 0.04 + 0.01\,\beta
\]
である。この期間の資産Aは点 $(\beta_A,\mu_A) = (0.375,\ 0.045)$ にプロット
される。$\beta = 0.375$ における証券市場線上の点は $(0.375,\ 0.04375)$ なので，
資産Aは証券市場線より\emph{上}にプロットされる。したがって，この期間の資産Aは
CAPMの理論値より高いリターンを上げており，\textbf{割安}であったと判断される。

\medskip
\textbf{ポイント。}(5)では資産AはCMLより下(非効率的な位置)にあり，(6)では
SMLより上(割安)にある。CMLは「効率的ポートフォリオ」だけが乗る線であるのに
対し，SMLは均衡時に「すべての資産」が乗る線であるという違いが，この問題の核心で
ある。
\end{answerbox}
